% Source: Graph_ML/Random_Walk_Similarity/random_walk_similarity.tex
\begin{notechapter}
% ---- chapter-local definitions ----
\makeatletter
\forgetenv{definition}\forgetenv{definition*}\expandafter\let\csname c@tcb@cnt@definition\endcsname\relax
\newtcbtheorem[number within=section]{definition}{Definition}{
  colback=blue!5!white,
  colframe=blue!75!black,
  fonttitle=\bfseries,
  breakable
}{random-walk-similarity:def}
\forgetenv{theorem}\forgetenv{theorem*}\expandafter\let\csname c@tcb@cnt@theorem\endcsname\relax
\newtcbtheorem[number within=section]{theorem}{Theorem}{
  colback=green!5!white,
  colframe=green!50!black,
  fonttitle=\bfseries,
  breakable
}{random-walk-similarity:thm}
\forgetenv{proposition}\forgetenv{proposition*}\expandafter\let\csname c@tcb@cnt@proposition\endcsname\relax
\newtcbtheorem[number within=section]{proposition}{Proposition}{
  colback=orange!5!white,
  colframe=orange!75!black,
  fonttitle=\bfseries,
  breakable
}{random-walk-similarity:prop}
\forgetenv{intuition}\newtcolorbox{intuition}{
  colback=yellow!10!white,
  colframe=yellow!50!black,
  title=Intuition,
  fonttitle=\bfseries,
  breakable
}
\forgetenv{keyidea}\newtcolorbox{keyidea}{
  colback=purple!5!white,
  colframe=purple!75!black,
  title=Key Idea,
  fonttitle=\bfseries,
  breakable
}
\forgetenv{example}\newtcolorbox{example}{
  colback=gray!5!white,
  colframe=gray!75!black,
  title=Example,
  fonttitle=\bfseries,
  breakable
}
\forgetcs\bA\newcommand{\bA}{\mathbf{A}}
\forgetcs\bD\newcommand{\bD}{\mathbf{D}}
\forgetcs\bP\newcommand{\bP}{\mathbf{P}}
\forgetcs\bI\newcommand{\bI}{\mathbf{I}}
\forgetcs\bq\newcommand{\bq}{\mathbf{q}}
\forgetcs\be\newcommand{\be}{\mathbf{e}}
\forgetcs\bS\newcommand{\bS}{\mathbf{S}}
\forgetcs\bM\newcommand{\bM}{\mathbf{M}}
\forgetcs\bT\newcommand{\bT}{\mathbf{T}}
\forgetcs\R\newcommand{\R}{\mathbb{R}}
\forgetcs\E\newcommand{\E}{\mathbb{E}}
\makeatother
% -----------------------------------
\chapter{Random-Walk Node Similarity}\label{ch:random-walk-similarity}
\chaptersubtitle{Following Hamilton's \emph{Graph Representation Learning}}

%=============================================================================
\section{Introduction}
%=============================================================================

In the study of graph-based similarity measures, we have seen methods based on counting paths between nodes, such as the Katz index and LHN similarity. These \emph{path-counting} methods use exact counts of walks between nodes to quantify similarity.

\textbf{Random walk methods} take a different approach: instead of counting all paths, they consider the \emph{probability} that a random walker starting at one node will reach another node. This probabilistic perspective offers several advantages:
\begin{itemize}
    \item Natural normalization (probabilities sum to 1)
    \item Interpretable as a stochastic process
    \item Connection to the famous PageRank algorithm
\end{itemize}

This document explains the \textbf{Personalized PageRank} approach to node similarity, as described in Hamilton's \emph{Graph Representation Learning}.

%=============================================================================
\section{Background: Random Walks on Graphs}
%=============================================================================

\subsection{The Stochastic Matrix}

\begin{definition}{Stochastic Matrix}{stochastic}
Given a graph $G = (V, E)$ with adjacency matrix $\bA$ and degree matrix $\bD = \text{diag}(d_1, \ldots, d_n)$, the \textbf{stochastic matrix} (or random walk transition matrix) is:
\begin{equation}
\bP = \bA \bD^{-1}
\label{random-walk-similarity:eq:stochastic}
\end{equation}
The entry $\bP[v, u] = \dfrac{\bA[v,u]}{d_u}$ gives the probability of transitioning from node $u$ to node $v$ in one step.
\end{definition}

\begin{intuition}
At each step, a random walker at node $u$ chooses uniformly at random among $u$'s neighbors. If $u$ has degree $d_u$, each neighbor has probability $1/d_u$ of being selected.

Note that $\bP$ is column-stochastic: each column sums to 1 (assuming no isolated nodes).
\end{intuition}

\subsection{Standard Random Walk Convergence}

For a connected, non-bipartite graph, if we let a random walk run indefinitely, the probability distribution over nodes converges to a \emph{stationary distribution}. This stationary distribution is proportional to node degrees---high-degree nodes are visited more often.

\begin{keyidea}
Without any modification, random walk probabilities converge to a \emph{global} measure (related to eigenvector centrality), not a node-specific similarity measure. To obtain pairwise similarity, we need the \textbf{Personalized PageRank} modification.
\end{keyidea}

%=============================================================================
\section{Personalized PageRank}
%=============================================================================

\subsection{The Core Idea: Random Walks with Restarts}

The key insight of Personalized PageRank is to periodically ``teleport'' the random walker back to a source node $u$. This creates a measure of how easily other nodes can be reached from $u$.

\subsubsection*{The Random Walk Process}

Consider a random walker that, at each discrete time step $t$:
\begin{itemize}
    \item With probability $c$: follows a random edge (standard random walk step)
    \item With probability $(1-c)$: ``teleports'' back to the source node $u$
\end{itemize}

Let $\bq_u^{(t)} \in \R^{|V|}$ be the probability distribution over nodes at time $t$, where $\bq_u^{(t)}[v]$ is the probability that the walker is at node $v$ at time $t$.

\subsubsection*{Deriving the Recurrence}

At time $t+1$, the walker is at node $v$ if either:
\begin{enumerate}
    \item The walker took a random step (probability $c$) and transitioned to $v$ from some neighbor $w$, or
    \item The walker teleported (probability $1-c$) and $v = u$.
\end{enumerate}

The probability of being at $v$ via a random step is:
\[
c \cdot \sum_{w \in V} \bP[v, w] \cdot \bq_u^{(t)}[w] = c \cdot (\bP \bq_u^{(t)})[v]
\]

The probability of being at $v$ via teleportation is:
\[
(1-c) \cdot \be_u[v] = \begin{cases} 1-c & \text{if } v = u \\ 0 & \text{otherwise} \end{cases}
\]

Combining these, the update rule is:
\begin{equation}
\bq_u^{(t+1)} = c \bP \bq_u^{(t)} + (1-c) \be_u
\label{random-walk-similarity:eq:ppr-dynamics}
\end{equation}

\subsubsection*{The Stationary Distribution}

We are interested in the \textbf{stationary distribution} $\bq_u$---the distribution that the walker converges to as $t \to \infty$. At stationarity, the distribution no longer changes:
\[
\bq_u^{(t+1)} = \bq_u^{(t)} = \bq_u
\]

Substituting into \cref{random-walk-similarity:eq:ppr-dynamics}:

\begin{definition}{Personalized PageRank}{ppr}
For a source node $u$, the \textbf{Personalized PageRank vector} $\bq_u \in \R^{|V|}$ is the stationary distribution satisfying:
\begin{equation}
\bq_u = c \bP \bq_u + (1 - c) \be_u
\label{random-walk-similarity:eq:ppr-recurrence}
\end{equation}
where:
\begin{itemize}
    \item $\be_u$ is the one-hot indicator vector for node $u$ (i.e., $\be_u[v] = 1$ if $v = u$, else 0)
    \item $c \in (0, 1)$ is the \textbf{damping factor} (typically $c \approx 0.85$)
    \item $\bq_u[v]$ gives the stationary probability that the random walker is at node $v$
\end{itemize}
\end{definition}

\begin{intuition}
The recurrence has a clear interpretation: the probability of being at any node $v$ is the sum of:
\begin{itemize}
    \item The probability of arriving via a random walk step: $c \cdot (\bP \bq_u)[v]$
    \item The probability of arriving via teleportation: $(1-c) \cdot \be_u[v]$
\end{itemize}
\end{intuition}

\subsection{Interpretation of the Damping Factor}

The parameter $c$ controls the trade-off between local and global information:

\begin{center}
\begin{tabular}{cl}
\toprule
\textbf{Value of $c$} & \textbf{Behavior} \\
\midrule
$c \to 0$ & Walker always restarts $\Rightarrow$ $\bq_u \to \be_u$ (only source node matters) \\
$c \to 1$ & Walker rarely restarts $\Rightarrow$ converges to global eigenvector centrality \\
$c \approx 0.85$ & Typical choice; balances local neighborhood with global structure \\
\bottomrule
\end{tabular}
\end{center}

\begin{intuition}
Think of $(1-c)$ as the probability of ``giving up'' on the current walk and starting over. With $c = 0.85$, the expected walk length before restart is $\frac{1}{1-c} = \frac{1}{0.15} \approx 6.67$ steps. See the box below for the derivation.
\end{intuition}

\begin{example}
\textbf{Expected walk length before restart.} At each step, the walker restarts with probability $(1-c)$ and continues with probability $c$. Let $L$ be the number of steps before a restart occurs. This is a \textbf{geometric random variable}: we're counting the number of ``successes'' (continuing) before the first ``failure'' (restarting).

The probability of taking exactly $k$ steps before restarting is:
\[
P(L = k) = c^k (1-c)
\]
This is because the walker must continue for $k$ steps (probability $c^k$) and then restart (probability $1-c$).

The expected value of a geometric distribution with ``failure'' probability $p = 1-c$ is:
\begin{align}
\E[L] &= \sum_{k=0}^{\infty} k \cdot c^k (1-c) = (1-c) \sum_{k=0}^{\infty} k \cdot c^k
\end{align}

To evaluate $\sum_{k=0}^{\infty} k \cdot c^k$, we use the identity\footnote{This follows from differentiating the geometric series $\sum_{k=0}^{\infty} c^k = \frac{1}{1-c}$ with respect to $c$.}:
\[
\sum_{k=0}^{\infty} k \cdot c^k = \frac{c}{(1-c)^2}
\]

Therefore:
\[
\E[L] = (1-c) \cdot \frac{c}{(1-c)^2} = \frac{c}{1-c}
\]

For $c = 0.85$: $\E[L] = \frac{0.85}{0.15} = 5.67$ steps on average before restart.

\textbf{Note:} There's an alternative convention where we count the restart step itself, giving $\E[L] = \frac{1}{1-c} = 6.67$. The document uses this convention.
\end{example}

%=============================================================================
\section{Closed-Form Solution}
%=============================================================================

\subsection{Deriving the Solution}

The recurrence in \cref{random-walk-similarity:eq:ppr-recurrence} can be solved analytically.

\begin{proposition}{Closed-Form Personalized PageRank}{ppr-closed}
The Personalized PageRank vector for source node $u$ is:
\begin{equation}
\bq_u = (1 - c)(\bI - c\bP)^{-1} \be_u
\label{random-walk-similarity:eq:ppr-closed}
\end{equation}
\end{proposition}

\begin{proof}
Starting from the recurrence $\bq_u = c \bP \bq_u + (1 - c) \be_u$, we rearrange:
\begin{align}
\bq_u - c \bP \bq_u &= (1 - c) \be_u \\
(\bI - c\bP) \bq_u &= (1 - c) \be_u \\
\bq_u &= (1 - c)(\bI - c\bP)^{-1} \be_u
\end{align}
This requires $(\bI - c\bP)$ to be invertible. See \cref{random-walk-similarity:app:invertibility} for the proof that this holds when $c < 1$.
\end{proof}

\subsection{Connection to Neumann Series}

Since $\|c\bP\| < 1$ in spectral norm (see \cref{random-walk-similarity:app:spectral-norm} for the definition), we can expand the inverse as a Neumann series (see \cref{random-walk-similarity:app:neumann}):
\begin{equation}
(\bI - c\bP)^{-1} = \sum_{k=0}^{\infty} (c\bP)^k = \bI + c\bP + c^2\bP^2 + c^3\bP^3 + \cdots
\label{random-walk-similarity:eq:neumann}
\end{equation}

\begin{intuition}
This expansion shows that Personalized PageRank is a weighted sum of random walk probabilities at different path lengths:
\begin{itemize}
    \item $\bP^0 = \bI$: staying at the source (weight 1)
    \item $\bP^1$: one-step transitions (weight $c$)
    \item $\bP^2$: two-step transitions (weight $c^2$)
    \item $\bP^k$: $k$-step transitions (weight $c^k$)
\end{itemize}
Longer paths contribute less due to the decaying weights $c^k$.
\end{intuition}

%=============================================================================
\section{Random Walk Similarity Measure}
%=============================================================================

\subsection{Definition}

Using Personalized PageRank, we can define a node-node similarity measure.

\begin{definition}{Random Walk Similarity}{rw-sim}
The \textbf{random walk similarity} between nodes $u$ and $v$ is:
\begin{equation}
\bS_{\text{RW}}[u, v] = \bq_u[v] + \bq_v[u]
\label{random-walk-similarity:eq:rw-sim}
\end{equation}
where $\bq_u[v]$ is the Personalized PageRank score of node $v$ with respect to source $u$.
\end{definition}

\begin{intuition}
The similarity between $u$ and $v$ is proportional to how likely we are to reach each node from a random walk (with restarts) starting from the other node. This is symmetric by construction.
\end{intuition}

\subsection{Why Symmetrize?}

Note that $\bq_u[v] \neq \bq_v[u]$ in general---the probability of reaching $v$ from $u$ may differ from reaching $u$ from $v$. The sum $\bq_u[v] + \bq_v[u]$ symmetrizes the measure, which is desirable for many applications where similarity should be a symmetric relation.

\begin{example}
Consider a directed social network where edges represent ``follows'' relationships. User $u$ might easily reach popular user $v$ (high $\bq_u[v]$), but $v$ might not easily reach $u$ (low $\bq_v[u]$). The symmetric measure captures both directions of reachability.
\end{example}

%=============================================================================
\section{Comparison with Path-Counting Methods}
%=============================================================================

\subsection{Random Walks vs. Katz Index}

\begin{center}
\begin{tabular}{lll}
\toprule
\textbf{Aspect} & \textbf{Katz Index} & \textbf{Random Walk (PPR)} \\
\midrule
Basic quantity & Path counts $\bA^k[u,v]$ & Transition probabilities $\bP^k[u,v]$ \\
Weighting & $\beta^k$ (decay factor) & $c^k$ (damping factor) \\
Matrix formula & $(\bI - \beta\bA)^{-1}$ & $(1-c)(\bI - c\bP)^{-1}$ \\
Degree bias & High (favors high-degree nodes) & Normalized (probabilities sum to 1) \\
Interpretation & Weighted path count & Stationary visit probability \\
\bottomrule
\end{tabular}
\end{center}

\begin{keyidea}
The key difference is that random walk methods use the \textbf{stochastic matrix} $\bP = \bA\bD^{-1}$ rather than the raw adjacency matrix $\bA$. This normalizes by degree, so high-degree nodes don't automatically dominate the similarity scores.
\end{keyidea}

\subsection{Connection to LHN}

Recall that LHN normalizes path counts by expected counts under a null model:
\[
\bS_{\text{LHN}}[u,v] \propto \frac{\bA^k[u,v]}{\E[\bA_{G'}^k[u,v]]}
\]

Random walk methods achieve a similar effect through a different mechanism: by using transition \emph{probabilities} rather than raw counts, the degree normalization is built into the random walk dynamics itself.

%=============================================================================
\section{Practical Considerations}
%=============================================================================

\subsection{Computational Complexity}

Computing $\bq_u$ for a single source node $u$ requires solving a linear system or performing power iteration. For all pairs of nodes:
\begin{itemize}
    \item \textbf{Exact computation}: $O(n^3)$ to compute $(\bI - c\bP)^{-1}$ (dense matrix inverse)
    \item \textbf{Iterative methods}: Power iteration converges quickly due to the spectral gap
    \item \textbf{Sparse approximations}: For large graphs, use truncated random walks or local push algorithms
\end{itemize}

\subsection{Choosing the Damping Factor}

The original PageRank paper used $c = 0.85$. Common choices:
\begin{itemize}
    \item $c = 0.85$: Standard choice, good balance of local/global
    \item $c = 0.5$: More emphasis on local neighborhood
    \item $c = 0.9$: More emphasis on global structure
\end{itemize}

The optimal choice depends on the application and graph structure.

%=============================================================================
\section{Summary}
%=============================================================================

\begin{keyidea}
\textbf{Random Walk Similarity} defines node similarity based on the stationary probabilities of random walks with restarts (Personalized PageRank).

\textbf{Key equations}:
\begin{align}
\text{Stochastic matrix:} \quad & \bP = \bA \bD^{-1} \\[0.5em]
\text{PPR recurrence:} \quad & \bq_u = c \bP \bq_u + (1 - c) \be_u \\[0.5em]
\text{Closed form:} \quad & \bq_u = (1 - c)(\bI - c\bP)^{-1} \be_u \\[0.5em]
\text{Similarity:} \quad & \bS_{\text{RW}}[u, v] = \bq_u[v] + \bq_v[u]
\end{align}

\textbf{Key insight}: Using transition probabilities ($\bP$) instead of raw adjacency ($\bA$) naturally normalizes for node degree, avoiding the degree bias present in methods like the Katz index.
\end{keyidea}

%=============================================================================
\chapterappendix
\section{Mathematical Background}
%=============================================================================

This appendix provides rigorous justification for the mathematical claims made in the main text.

\subsection{Spectral Norm}
\label{random-walk-similarity:app:spectral-norm}

\begin{definition}{Spectral Norm}{spectral-norm}
The \textbf{spectral norm} (or operator 2-norm) of a matrix $\bA \in \R^{n \times n}$ is:
\[
\|\bA\| = \max_{\mathbf{x} \neq \mathbf{0}} \frac{\|\bA \mathbf{x}\|_2}{\|\mathbf{x}\|_2} = \sigma_{\max}(\bA)
\]
where $\sigma_{\max}(\bA)$ is the largest singular value of $\bA$.
\end{definition}

For a square matrix, the spectral norm equals the largest singular value. For a symmetric matrix, this equals the largest absolute eigenvalue: $\|\bA\| = \max_i |\lambda_i|$.

\begin{proposition}{Spectral Norm of Stochastic Matrix}{stochastic-norm}
If $\bP$ is a column-stochastic matrix (each column sums to 1), then $\|\bP\| \leq 1$.
\end{proposition}

\begin{proof}
For any vector $\mathbf{x}$, consider $\|\bP\mathbf{x}\|_1$. Since each column of $\bP$ sums to 1 and has non-negative entries:
\[
\|\bP\mathbf{x}\|_1 = \sum_i |(\bP\mathbf{x})_i| \leq \sum_i \sum_j \bP_{ij} |x_j| = \sum_j |x_j| \sum_i \bP_{ij} = \sum_j |x_j| = \|\mathbf{x}\|_1
\]
Thus $\|\bP\|_1 \leq 1$. Similarly, $\|\bP^\top\|_1 \leq 1$ (row sums are at most 1 for column-stochastic matrices of undirected graphs). Since $\|\bP\|^2 \leq \|\bP\|_1 \|\bP^\top\|_1$, we have $\|\bP\| \leq 1$.
\end{proof}

\subsection{Eigenvalues of Stochastic Matrices}
\label{random-walk-similarity:app:eigenvalues}

\begin{proposition}{Eigenvalue Bound for Stochastic Matrices}{eigenvalue-bound}
If $\bP$ is a column-stochastic matrix, then all eigenvalues $\lambda$ of $\bP$ satisfy $|\lambda| \leq 1$.
\end{proposition}

\begin{proof}
Let $\lambda$ be an eigenvalue of $\bP$ with eigenvector $\mathbf{v}$. Then $\bP\mathbf{v} = \lambda \mathbf{v}$.

By the Gershgorin circle theorem, every eigenvalue lies in at least one Gershgorin disc. For column $j$, the disc is centered at $\bP_{jj}$ with radius $\sum_{i \neq j} \bP_{ij}$. Since $\bP$ is column-stochastic:
\[
\bP_{jj} + \sum_{i \neq j} \bP_{ij} = 1
\]
So each disc is contained in the closed unit disc $\{z \in \mathbb{C} : |z| \leq 1\}$. Therefore $|\lambda| \leq 1$.
\end{proof}

\subsection{Invertibility of $(\bI - c\bP)$}
\label{random-walk-similarity:app:invertibility}

\begin{proposition}{Invertibility}{invertibility}
For $c \in (0, 1)$ and $\bP$ a column-stochastic matrix, the matrix $(\bI - c\bP)$ is invertible.
\end{proposition}

\begin{proof}
The matrix $(\bI - c\bP)$ is invertible if and only if $0$ is not an eigenvalue, i.e., if and only if $\det(\bI - c\bP) \neq 0$.

The eigenvalues of $(\bI - c\bP)$ are $\{1 - c\lambda : \lambda \text{ is an eigenvalue of } \bP\}$.

By \cref{random-walk-similarity:prop:eigenvalue-bound}, $|\lambda| \leq 1$ for all eigenvalues of $\bP$. Therefore:
\[
|1 - c\lambda| \geq 1 - c|\lambda| \geq 1 - c > 0
\]
since $c < 1$. Thus no eigenvalue of $(\bI - c\bP)$ is zero, so the matrix is invertible.
\end{proof}

\subsection{Neumann Series}
\label{random-walk-similarity:app:neumann}

\begin{theorem}{Neumann Series}{neumann}
If $\bM$ is a square matrix with $\|\bM\| < 1$ (in any submultiplicative norm), then $(\bI - \bM)$ is invertible and:
\begin{equation}
(\bI - \bM)^{-1} = \sum_{k=0}^{\infty} \bM^k = \bI + \bM + \bM^2 + \bM^3 + \cdots
\end{equation}
The series converges absolutely.
\end{theorem}

\begin{proof}
\textbf{Step 1: Convergence.} Since $\|\bM\| < 1$ and the norm is submultiplicative, $\|\bM^k\| \leq \|\bM\|^k$. The series $\sum_{k=0}^{\infty} \|\bM\|^k$ is a geometric series with ratio $< 1$, so it converges. By comparison, $\sum_{k=0}^{\infty} \bM^k$ converges absolutely.

\textbf{Step 2: Verification.} Let $\bS = \sum_{k=0}^{\infty} \bM^k$. Then:
\begin{align}
(\bI - \bM) \bS &= (\bI - \bM) \sum_{k=0}^{\infty} \bM^k = \sum_{k=0}^{\infty} \bM^k - \sum_{k=0}^{\infty} \bM^{k+1} \\
&= \sum_{k=0}^{\infty} \bM^k - \sum_{k=1}^{\infty} \bM^k = \bM^0 = \bI
\end{align}
Similarly, $\bS(\bI - \bM) = \bI$. Therefore $\bS = (\bI - \bM)^{-1}$.
\end{proof}

\begin{intuition}
The Neumann series is the matrix analogue of the geometric series $\frac{1}{1-x} = 1 + x + x^2 + \cdots$ for $|x| < 1$.

Applying this to $\bM = c\bP$ where $\|c\bP\| = c\|\bP\| \leq c < 1$:
\[
(\bI - c\bP)^{-1} = \sum_{k=0}^{\infty} c^k \bP^k
\]
This shows that Personalized PageRank weights $k$-step transition probabilities by $c^k$.
\end{intuition}

\subsection{Convergence of Power Iteration}
\label{random-walk-similarity:app:power-iteration}

The recurrence $\bq_u^{(t+1)} = c\bP\bq_u^{(t)} + (1-c)\be_u$ converges to the stationary distribution.

\begin{proposition}{Convergence}{convergence}
For any initial distribution $\bq_u^{(0)}$, the sequence $\bq_u^{(t)}$ converges to $\bq_u = (1-c)(\bI - c\bP)^{-1}\be_u$.
\end{proposition}

\begin{proof}
Let $\bT: \mathbf{x} \mapsto c\bP\mathbf{x} + (1-c)\be_u$ be the update operator. The fixed point satisfies $\bq_u = \bT(\bq_u)$.

The error at step $t$ is:
\begin{align}
\bq_u^{(t)} - \bq_u &= \bT(\bq_u^{(t-1)}) - \bT(\bq_u) \\
&= c\bP(\bq_u^{(t-1)} - \bq_u)
\end{align}

By induction: $\bq_u^{(t)} - \bq_u = (c\bP)^t (\bq_u^{(0)} - \bq_u)$.

Since $\|c\bP\| \leq c < 1$:
\[
\|\bq_u^{(t)} - \bq_u\| \leq \|c\bP\|^t \|\bq_u^{(0)} - \bq_u\| \leq c^t \|\bq_u^{(0)} - \bq_u\| \to 0
\]
as $t \to \infty$. The convergence is geometric with rate $c$.
\end{proof}

%=============================================================================
\section{References}
%=============================================================================

\begin{itemize}
    \item Hamilton, W. L. (2020). \emph{Graph Representation Learning}. Morgan \& Claypool.
    \item Page, L., Brin, S., Motwani, R., \& Winograd, T. (1999). The PageRank citation ranking: Bringing order to the web. Stanford InfoLab.
    \item Leskovec, J., Rajaraman, A., \& Ullman, J. D. (2020). \emph{Mining of Massive Datasets}. Cambridge University Press.
\end{itemize}
\end{notechapter}
